%%%%%%%%%%%%%%%%%%%%%%%%%%%%%%%%%%%%%%%%%%%%%%%%%%%%%%
\newpage
\pagestyle{fancy}

\chapter{Title of the Chapter 2}
\minitoc
\newpage
\label{chap.2}

\section{Flow Diverter Stents}\label{chap2:sec1}

Flow Diverter stents (FDs) have become a viable alternative to coiling of IA, especially in cases of giant wide-necked and fusiform aneurysms. The use of coils in those particular cases bears a higher risk of coil compaction or detachment. FDs, on the contrary, gain their stability through the aposition with the parent artery. The operating doctor guides the microcatheter towards a position downstream of the IA. The distal end of the FD is then pushed and expanded, ideally making contact with the complete vessel cross-section. By carefully withdrawing the catheter and pushing the FD, the device is then deployed across the aneurysm. One important criterion for a successful deployment is achieving an optimal wall aposition, i.e.\, setting a stent whose radial force is sufficiently large to ensure a firm anchoring in the artery. Neuroradiologists can control the aposition both with the sizing of the device, as well as with the force exerted during deployment. A compact segment of the stent naturally offers a reduced porosity, which is also exploited if a higher exclusion of the flow is sought at the aneurysm neck. According to the testimony of neuroradiologists and medical students, the 

\newpage
\bibliographystyle{elsarticle-num}
\bibliography{chapter2/Biblio}
